\documentclass[11pt,letterpaper]{article}
\usepackage[T1]{fontenc}
\usepackage[utf8]{inputenc}
\usepackage{lmodern}
\usepackage[margin=1in]{geometry}
\usepackage{booktabs}
\usepackage{longtable}
\usepackage{enumitem}
\usepackage{listings}
\usepackage{xcolor}
\usepackage[hidelinks]{hyperref}

\definecolor{backcolour}{rgb}{0.96,0.96,0.94}
\definecolor{cautionred}{rgb}{0.55,0,0}
\definecolor{okgreen}{rgb}{0,0.45,0}

% See BASIC_CAN_COMMANDS.tex for why these options are set: everything in a
% listing must survive being selected from the PDF and pasted into a terminal.
\lstset{
    basicstyle=\ttfamily\small,
    backgroundcolor=\color{backcolour},
    numbers=none,
    frame=single,
    framesep=5pt,
    columns=fullflexible,
    keepspaces=true,
    upquote=true,
    showstringspaces=false,
    breaklines=false,
    xleftmargin=0pt,
    aboveskip=8pt,
    belowskip=8pt
}

\newenvironment{caution}[1]{%
    \par\addvspace{8pt}%
    \noindent{\color{cautionred}\rule{\linewidth}{0.8pt}}\par\vspace{2pt}
    \noindent\textbf{\color{cautionred}Caution: #1}\par\vspace{2pt}
}{%
    \par\vspace{2pt}\noindent{\color{cautionred}\rule{\linewidth}{0.4pt}}%
    \par\addvspace{8pt}%
}
\newcommand{\passif}[1]{\par\vspace{2pt}\noindent\textbf{\color{okgreen}Pass:} #1\par\vspace{4pt}}
\newcommand{\failif}[1]{\noindent\textbf{If it fails:} #1\par\vspace{6pt}}

\title{NEMA 17 / MKS SERVO42D --- Bring-Up SOP \\ \large Arctos Robotic Arm --- joints A, B and C}
\date{September 2026}

\begin{document}
\maketitle

Follow these steps in order to take a NEMA 17 joint from powered-off to trusted for
motion. Every step has a pass condition; if one fails, stop there rather than
continuing.

Background, protocol details and troubleshooting live in
\texttt{Arctos\_CAN\_Interface.pdf}. Bare terminal commands live in
\texttt{BASIC\_CAN\_COMMANDS.pdf}.

\section*{Which joint you have}

\begin{center}
\begin{tabular}{llll}
\toprule
\textbf{Joint} & \textbf{CAN ID} & \textbf{Gear ratio} & \textbf{Role} \\ \midrule
A & \texttt{0x04} & 48.0:1  & Elbow \\
B & \texttt{0x05} & 67.82:1 & Wrist rotation \\
C & \texttt{0x06} & 67.82:1 & Wrist joint \\
\bottomrule
\end{tabular}
\end{center}

\begin{caution}{Confirm the ID by scan, and reprogram the panel if it disagrees}
The table above is the intended assignment. A controller's panel can hold a
different value: on 2026-09-02 Joint B was found answering on \texttt{0x06}, which
is Joint C's ID, because that driver's panel had been saved wrong.

If the scan in Step 3 disagrees with this table, \textbf{go back to Step 1 and
reprogram the panel} rather than carrying on against the wrong ID. Two joints
sharing an ID will collide, and a frame sent to the wrong ID drives the wrong
joint.
\end{caution}

Wherever a command below shows \texttt{0x06} and \texttt{67.82}, substitute your own
joint's CAN ID and gear ratio.

\section*{Step 0 --- Before applying power}
\begin{enumerate}[leftmargin=*]
    \item Confirm the controller is the \textbf{CAN variant} of the MKS SERVO42D.
          Read the part number on the board. The RS485 variant has no CAN
          transceiver and will never respond, whatever you do in software.
    \item Everything powered off, measure CAN\_H to CAN\_L. Expect about
          60~$\Omega$.
    \item Confirm only CAN\_H, CAN\_L and GND run between the CANable and the
          controller. Do not connect the CANable's power lines.
    \item Confirm the joint is mechanically clear, and that you know its travel
          limits if it is mounted on the arm.
\end{enumerate}
\passif{About 60~$\Omega$ across the bus, and a confirmed CAN-variant board.}
\failif{120~$\Omega$ means only one terminator is present. An open circuit means
broken wiring. Fix it before continuing.}

\section*{Step 1 --- Configure the controller panel}

Power the controller and set these on its own display, then \textbf{save}:

\begin{center}
\begin{tabular}{ll}
\toprule
CAN ID & the ID for this joint \\
CAN Rate & \texttt{500} \\
CAN Response, shown as CANRSP & \texttt{ENABLE} \\
CAN CRC & \texttt{SUM-8} \\
Working current & 1.2~A or below, the NEMA 17's rating \\
Working mode & note what it is set to \\
\bottomrule
\end{tabular}
\end{center}

Power-cycle the controller after saving. The CAN ID in particular has to be
committed, not merely displayed.

\begin{caution}{Set the working current deliberately}
A controller in this project was found storing a 3000~mA default against a motor
rated 1.2~A, roughly two and a half times over. Check this on every new controller.
\end{caution}

\begin{caution}{A working-mode mismatch imitates a communication fault}
If the panel's working mode disagrees with what a motion command assumes, the
controller will accept a frame and acknowledge its checksum without turning the
shaft. Confirm the mode before concluding that a ``no motion'' result is a wiring or
CAN problem.
\end{caution}
\passif{Settings survive the power cycle.}

\section*{Step 2 --- Bring up the CAN interface}

\begin{lstlisting}
lsusb
ls -l /dev/serial/by-id/
sudo pkill slcand
sudo ip link delete can0 2>/dev/null
sudo slcand -o -c -s6 /dev/serial/by-id/usb-Openlight_Labs_CANable2* can0
sudo ip link set can0 txqueuelen 1000
sudo ip link set up can0
ip -br link show can0
\end{lstlisting}

\passif{The last command prints \texttt{can0} followed by \texttt{UP}.}
\failif{Re-check the device name; it changes whenever the adapter re-enumerates.
\texttt{bitrate 0} in \texttt{ip -details} output is normal on slcan and is not a
fault.}

\section*{Step 3 --- Confirm the real CAN ID}

Never assume the ID from a config file or a table.

\begin{lstlisting}
cd ~/arctos_ws/src/arctos_can_control/arctos_can_control
python3 can_id_scan.py --channel can0 --ids 1-16
\end{lstlisting}

This sends only the read-only \texttt{0x31} query. It never enables coils and never
commands motion, so it is safe at any time.

\passif{Exactly one responder, and its ID is the one you set in Step 1.}
\failif{No responders: re-check the panel settings, the wiring, and the board's part
number. Several responders: another driver is powered on the bus with a colliding
ID. Resolve that before going further.}

\section*{Step 4 --- Read-only health check}

With the ID confirmed, check the link is trustworthy before commanding anything.

\begin{lstlisting}
python3 joint_reliability_test.py --channel can0 --can-id 0x06 \
    --gear-ratio 67.82 --samples 200
\end{lstlisting}

With the joint at rest, expect:

\begin{center}
\begin{tabular}{lll}
\toprule
\textbf{Query} & \textbf{Opcode} & \textbf{Expected} \\ \midrule
Encoder          & \texttt{0x31} & stable; zero jitter, zero drift \\
Velocity         & \texttt{0x32} & zero \\
Position error   & \texttt{0x39} & small, tens of counts \\
Enable state     & \texttt{0x3A} & \texttt{01}; these drivers power up enabled \\
Shaft protection & \texttt{0x3E} & \texttt{00} \\
\bottomrule
\end{tabular}
\end{center}

\passif{100\% response rate, zero checksum failures, zero encoder jitter at rest.}
\failif{Any checksum failure or dropped response is a wiring or bitrate problem. Do
not proceed to motion.}

\section*{Step 5 --- Check the mechanism by hand}

Do this before the motor ever drives the joint. It separates mechanical faults from
electrical ones, which the bus on its own cannot distinguish.

\begin{enumerate}[leftmargin=*]
    \item Disable the coils. For \texttt{0x06}:
\begin{lstlisting}
cansend can0 006#F300F9
\end{lstlisting}
    \item Confirm the state is now \texttt{00}:
\begin{lstlisting}
cansend can0 006#3A40
\end{lstlisting}
    \item Turn the joint output by hand, gently, both directions, watching the
          encoder.
\end{enumerate}

\begin{caution}{Only drop coils if the joint is not gravity-loaded}
On an assembled arm this removes holding torque and the joint may fall. Support it,
or skip this step if it cannot be done safely.
\end{caution}

\passif{The encoder tracks your hand in both directions and the motion feels smooth.
This proves the gearbox transmits and that nothing is skipping.}
\failif{Gritty or catching resistance means the gear mesh needs a plastic-safe PTFE
or silicone grease. No encoder movement while the output turns means the motor shaft
is not following: check the coupling, the grub screw and the bore. A hard stop at
one repeatable angle is usually a misassembled gear rather than roughness --- on
Joint B a gear seated in the wrong position produced a hard limit at exactly the
same angle on eight consecutive attempts.}

\section*{Step 6 --- First commanded motion}

Paste your joint's emergency stop into a second terminal, ready but unsent, before
you send anything. For \texttt{0x06}:

\begin{lstlisting}
cansend can0 006#F7FD
\end{lstlisting}

Re-enable the coils, then command one small move:

\begin{lstlisting}
python3 joint_unstick_move_test.py --can-id 0x06 --gear-ratio 67.82 \
    --step-degrees 0.25 --num-steps 1 --max-unstick 0 \
    --band-lo -1 --band-hi 1
\end{lstlisting}

\begin{caution}{Always set the band explicitly}
Some joints have a \texttt{home\_position} and an \texttt{opposite\_limit} of
\texttt{0.0} in \texttt{arctos\_controller.yaml}, which produces an inverted and
unusable safe band. Pass the band on the command line yourself, every time.
\end{caution}

\passif{The step reports about 100\% of commanded and the status frames show
\texttt{FD=[1, 2]}. Both frames matter: \texttt{FD 01} on its own means the command
started but never finished.}
\failif{Do not retry blindly --- that only heats the motor. Go back to Step 5.}

\section*{Step 7 --- Stepped travel test}

Extend to real travel, one direction and then the other. Keep the step size and
count inside the joint's safe range.

\begin{lstlisting}
python3 joint_unstick_move_test.py --can-id 0x06 --gear-ratio 67.82 \
    --step-degrees 2.0 --num-steps 9 --max-unstick 0 \
    --band-lo -25 --band-hi 25 --speed 25
\end{lstlisting}

Then the same distance back, to close the loop:

\begin{lstlisting}
python3 joint_unstick_move_test.py --can-id 0x06 --gear-ratio 67.82 \
    --step-degrees -2.0 --num-steps 9 --max-unstick 0 \
    --band-lo -25 --band-hi 25 --speed 25
\end{lstlisting}

\passif{Every step 99--101\% of commanded with \texttt{FD 02}, and the round trip
returns to the starting position within a few encoder counts.}
\failif{A step that under-delivers followed by a step that delivers nothing means
the joint is against something, or the motor has lost drive. Stop and diagnose. Do
not increase torque and do not retry.}

\section*{Step 8 --- Sweep test, if the mesh is rough}

Only needed if Step 7 showed a rough patch that a hand could push through. This
sweeps back and forth across the spot for a fixed time, backing off and re-attempting
whenever a step under-delivers.

\begin{lstlisting}
python3 joint_sweep_unstick_test.py --can-id 0x06 --gear-ratio 67.82 \
    --sweep-lo -10 --sweep-hi 10 --step-degrees 1.0 \
    --backoff-degrees 0.5 --band-lo -25 --band-hi 25 --duration-s 150
\end{lstlisting}

\passif{Zero unstick cycles used, and every step 99--101\%.}
\failif{If the joint stops at the \emph{same} angle every time, that is a mechanical
limit, not roughness --- open the gearbox rather than working at it. Repeated
back-and-forth through a hard particle embeds it and wears printed teeth.}

\section*{Step 9 --- Measure the speed ceiling}

\begin{lstlisting}
python3 joint_speed_limit_test.py --can-id 0x06 --gear-ratio 67.82 \
    --arc-lo -24 --arc-hi 24 --band-lo -25 --band-hi 25 \
    --speeds 100,150,175,200,250 --accel 8 --pause-s 15
\end{lstlisting}

The usable limit is the highest commanded speed where actual rpm still tracks it to
within about 1\%.

\begin{caution}{Keep the duty cycle down}
Running a long ladder of trials back to back overheats the motor. Fifteen seconds
between trials is reasonable; one second is not. Stop as soon as you have the
answer, and put a hand on the motor part way through --- there is no temperature
opcode, so nothing in software will warn you.
\end{caution}

Measured result for Joint B, through its 67.82:1 gearbox:

\begin{center}
\begin{tabular}{ll}
\toprule
Tracks commanded within 1\% up to & 150 rpm \\
Roll-off begins                   & between 150 and 165 rpm \\
Saturation ceiling                & about 165 rpm \\
\bottomrule
\end{tabular}
\end{center}

A bare Servo42D on the bench tracks to 225 rpm and saturates near 310. A geared
joint saturates well below that, so measure each joint rather than inheriting the
bench figure.

\section*{Step 10 --- Acceptance}

The joint is ready for use when all of these hold:

\begin{center}
\begin{tabular}{ll}
\toprule
\textbf{Check} & \textbf{Criterion} \\ \midrule
CAN ID            & confirmed by scan \\
Comms             & 100\% response, zero checksum failures \\
Encoder at rest   & zero jitter, zero drift \\
Hand check        & smooth, encoder tracks both directions \\
Per-step accuracy & 99--101\% of commanded \\
Status frames     & \texttt{FD 02} on every step \\
Round trip        & returns to start within a few counts \\
Shaft protection  & still \texttt{00} \\
Speed limit       & measured and written down \\
\bottomrule
\end{tabular}
\end{center}

\section*{Step 11 --- Operating limits}

\begin{itemize}[leftmargin=*]
    \item \textbf{Stay at or below the measured saturation speed.} Commanding more
          buys no extra speed and only adds current and heat.
    \item \textbf{Use \texttt{0xFD} for all motion.} Do not use \texttt{0xF5}.
    \item \textbf{Never retry into a stall.} Back off and re-attempt once, or stop.
    \item \textbf{Disable coils between tests}, unless the joint is gravity-loaded,
          in which case leave it holding.
    \item \textbf{Never enable shaft protection} (\texttt{0x88}). It has a history of
          false positives and latches the driver.
\end{itemize}

\section*{Step 12 --- Shutdown}

Stop any running script with Ctrl+C first.

\begin{lstlisting}
cansend can0 006#F300F9
sudo ip link set can0 down
sudo pkill slcand
\end{lstlisting}

Then unplug the adapter. Skip the first line if the joint is gravity-loaded and
needs to keep holding.

\end{document}
